\chapter{Progetti Filtri IIR}

\begin{figure}[h]
    \centering
    \includegraphics[scale = 0.8]{schema-di-realizzazione-di-un-filtro-IIR-nella-forma-Diretta-I.jpg}
\end{figure}  

\newpage 

\section{Progetto di filtri}

\footnote{ 
    Slide del prof | 9 - Progetto Filtri IIR 1 | pag 2 - 3\\
    Appunti | 2025-11-11 | pag 4   
}

Come scritto nelle pagine precedenti, in questo corso andremo affrontare i filtri discreti, che, essendo sistemi (tristi) lineari (o meglio LTI), 
sono selettivi in frequenza, capaci di rimuovere e non aggiungere alcune componenti in frequenza. \newline 

Solo i sistemi non lineari posso aggiungere componenti in frequenza, ma non verrano affrontati in questo corso. \newline 

Come ogni progetto, per la progettazione dei filtri, bisogna rispettare delle specifiche. \newline 

\begin{tcolorbox}
Come ribadito dal prof a lezione, queste specifiche possono essere date da un cliente o da uno sceicco con tanti soldi come il seguente: \newline 


\includegraphics[scale = 0.6]{Ciccio gamer89 lo sceicco.png}

Magari servirà un filtro LTI per filtrare la voce dello sceicco da questa canzone: \newline 

YTP Music - CiccioTrapper89 - HakunamaTATATATATATA! by Frederick Frankenstein\\
\url{https://www.youtube.com/watch?v=l6X7bur8MsQ&list=RDl6X7bur8MsQ&start_radio=1} \newline 

Oppure in questa fantastico suo discorso: \newline 

un po' di "Hakuna Matata Ragazzi" by UnusualValve\\
\url{https://www.youtube.com/watch?v=S7d9XmH-IYU} \newline 

(Questi appunti stanno diventando l'incontro tra Wikipedia e Nonciclopedia)
\end{tcolorbox}

Tipicamente, le specifiche vengono date in analogico e sono le seguenti: 

\begin{itemize}
    \item si specifica un determinato andamento per la risposta in modulo cioè il modulo del filtro LTI $\abs{H \left(e^{\jmath \omega}\right)}$ deve essere di un certo valore per pulsazioni $\abs{\omega} < \pi$ (perchè consideriamo lo spettro del segnale nelle pulsazioni $\omega$ compreso nell'intervallo $\left[ - \pi, + \pi\right]$) 
    \item si specificano i vincoli sulla fase e sul ritardo di gruppo (ad esempio fase lineare, fase "zero", fase minima, ecc...)
    \item si specificano i vincoli sulla precisione sul rapporto segnale rumore (scritto anche come SNR Signal Noise Ratio o S/R) o sulla aritmetica da utilizzare 
    \item si specificano i vincoli tecnologici (cioè la complessità, la velocità, il costo economico della soluzione, ecc...)
\end{itemize}

\newpage 

\begin{tcolorbox}
    Il terzo e quarto punto, in questo corso non verranno affrontati perchè sono degli argomenti in cui è necessario uno studio approfondito: 
    ci concentreremo sui primi due. \newline 

    Il rapporto segnale rumore può essere anche definito come errore di quantizzazione. 
\end{tcolorbox}

Per quanto riguarda la progettazione del filtro digitale, cioè $H \left(z\right)$, bisognerà svolgere queste scelte: 

\begin{tcolorbox}[colback=red!5!white,colframe=red!75!black] \label{q: 1}
Giusta la frase scritta? 
\end{tcolorbox}

\begin{itemize}
    \item si sceglie se una realizzazione FIR, Finite Impulse Response o risposta impulsiva finita, o una realizzazione IIR, Infinite Impulse Response o risposta impulsiva infinita 
    \begin{tcolorbox}
        Come ribadito in \ref{nota: differenza tra FIR e IIR} \nameref{nota: differenza tra FIR e IIR} pag \pageref{nota: differenza tra FIR e IIR}, 
        i FIR hanno tanti coefficienti $b_k$ ma garantiscono sempre stabilità del filtro, a differenza degli IIR che hanno, in confronto ai FIR, meno coefficienti $a_k$ e $b_k$ ma deve essere garantita la stabilità. 
    \end{tcolorbox}

    \item si calcola la relativa funzione $H \left(z\right)$
\end{itemize}

A livello del progetto di struttura, si sceglie la struttura più adatta in basa ai vincoli, cioè l'implementazione del filtro. \newline 

Alla fine, si svolge una verifica delle caratteristiche ottenute: in caso contrario, si ritorna ai passaggi precedenti affinche le specifiche non vengono soddisfatte. \newline 

\newpage 

\subsection{Specifiche "in frequenza"} \label{subsec: Specifiche "in frequenza"}

\footnote{ 
    Slide del prof | 9 - Progetto Filtri IIR 1 | pag 4\\
    Appunti | 2025-11-11 | pag 4   
} 

Molto spesso, le specifiche vengono date in analogico, nel mondo reale, e non in digitale. \newline 

In particolare, le pulsazioni vengono date "in frequenza" in $\omega$. \newline 

\begin{tcolorbox}
    Tecnicamente è sbagliato dire che le pulsazioni sono uguali alla frequenza, ma ricordati che la relazione tra $\omega$ pulsazione e f frequenza è la seguente: 

    {
        \Large 
        \begin{equation}
            \omega = 2 \pi f
        \end{equation}
    }

    quindi, c'è di mezzo solo un fattore moltiplicativo di $2 \pi$
\end{tcolorbox}

Considerando la seguente figura di un filtro passa-basso prototipo con le specifche, 
con risposta normalizzata nella banda passante, cioè $\abs{H \left(e^{\jmath \omega}\right)} = 1$: 

\begin{figure}[h]
    \centering
    \includegraphics[scale = 1]{prototipo filtro passa basso.png}
\end{figure}  

possiamo svolgere diverse considerazioni. \newline 

\begin{tcolorbox}
    Vengono date tutte queste specifiche per un filtro passa-basso perchè i filtri passa-alto e passa-banda si riportano al filtro passa-basso. 
\end{tcolorbox}

Idealmente vorremmo avere un filtro in cui la banda di transizione fosse nulla perchè vorremmo avere un filtro rettangolare in frequenza, 
ma ciò non è possibile perchè, ricordando la relazione tra il dominio del tempo e Fourier, se si ha una transizione brusca in Fourier si avrà un andamento nel tempo infinito. \newline 

\begin{tcolorbox}[colback=red!5!white,colframe=red!75!black] \label{q: 2}
Giusta la frase scritta? 
\end{tcolorbox}

Per questo motivo, nel filtro prototipo, quindi anche nelle specifiche, è necessario avere una banda di transizione non nulla (nella figura, quelle pulsazioni $\omega$ comprese nell'intervallo $\left[ \omega_p, \omega_s\right]$). \newline 

\begin{tcolorbox}[colback=red!5!white,colframe=red!75!black] \label{q: 3}
C'è un motivo per cui il modulo del filtro nella banda di transizione, almeno nella figura, è lineare? Ha a che fare con la fase?  
\end{tcolorbox}

\newpage

Essendoci una banda di transizione, la risposta in modulo del filtro, cioè $\abs{H \left(e^{\jmath \omega}\right)}$, non viene asseganata come funzione di 
$\omega$ per $\abs{\omega} < \pi$, ma viene assegnata per punti e/o attraverso dei seguenti vincoli, che devono essere rispettati: 

{
    \Large 
    \begin{equation}
        \begin{cases}
            \left(1 - \delta_1\right) \leq \abs{H \left(e^{\jmath \omega}\right)} \leq \left(1 - \delta_1\right) 
            \text{ per }
            \abs{\omega} \leq \pi 
            \\
            \abs{H \left(e^{\jmath \omega}\right)} \leq \delta_2 
            \text{ per }
            \omega_s \leq \abs{\omega} \leq \pi  
        \end{cases}
    \end{equation}
}

dove: 

\begin{itemize}
    \item $\delta_1$ e $\delta_2$ devono essere il più piccoli possibile, in modo che il filtro tenda all'andamento di un filtro rettangolare ideale 
    \item $\abs{H \left(e^{\jmath \omega}\right)}$, oltre ad essere compreso $\left(1 - \delta_1\right)$ e $\left(1 + \delta_1\right)$, può avere qualsiasi andamento in modulo ma, ricordando di mantenere, se possibile, la linearità della fase 
    \item in banda oscura, l'andamento $\abs{H \left(e^{\jmath \omega}\right)}$ può essere qualsiasi sia nei moduli che nella fase, l'importante il modulo sia minore di $\delta_2$
\end{itemize}

\newpage 

\section{Progetto di filtri IIR}

\footnote{ 
    Slide del prof | 9 - Progetto Filtri IIR 1 | pag 5\\
    Appunti | 2025-11-11 | pag 4   
} 

Per la progettazione di un filtro IIR, generalmente, si applicano una tra le due tecniche di progettazione: 

\begin{itemize}
    \item tecniche di trasformazione (le più utilizzate per i filtri IIR)
    \item tecniche dirette (le più utilizzate per i filtri FIR)
\end{itemize}

Per tecniche di trasformazione si intende che, inizialmente, si progetta il filtro (o un suo prototipo) nel dominio di Laplce, cioè nel dominio s, 
cioè si progetta un filtro analogico corrispondente. \newline 

Poi, si trasforma il filtro analogico, cioè $H_c \left(s\right)$ del filtro continuo in una funzione $H \left(z\right)$, cioè nel dominio discreto z, che è proprio il filtro voluto in digitale. \newline 

Il vantaggio maggiore delle tecnice di trasformazione è che, grazie alla teoria della approssimazione, si dispone di forme chiuse per il filtro analogico $H_c \left(s\right)$ (cosa che non è possibile fare nel dominio z a differenza del dominio s). \newline 

\begin{tcolorbox}[colback=red!5!white,colframe=red!75!black] \label{q: 4}
Cosa si intende per forme chiuse in s ? \newline 

Da Wikipedia: \newline 

In matematica, per forma chiusa può intendersi:
\begin{itemize}
    \item una forma differenziale chiusa, cioè tale che $d\omega = 0$ 
    \item un'espressione chiusa, cioè definita da un numero finito di applicazioni di certe funzioni, intesa come antitesi di forma ricorsiva.  
\end{itemize}

Cioè, per quanto riguarda la progettazione dei filtri IIR e le tecniche di trasformazione, per forme chiuse si intende che non possiamo rappresentare tutti i filtri analogici $H_c \left(s\right)$ dal dominio s al dominio z, ma solo alcuni? 

\end{tcolorbox}

Lo svantaggio maggiore delle tenciche di trasformazione è che si ha uno scarso controllo delle prestazione ottenuto, quindi è difficile ottimizzare $H \left(z\right)$ da $H_c \left(s\right)$. \newline 

Invece per tecniche dirette, significa che si progetta il filtro direttamente nel dominio z, ponendo dei vincoli alla funzione $H \left( z\right)$. \newline 

In altre parole, è possibile gestire i coefficienti della funzione $H \left(z\right)$. \newline 

I vantaggi delle tecniche dirette sono che si ha un buon controllo delle prestazioni e si ha hanno delle prestazioni. \newline 

Gli svantaggi delle tecniche dirette sono che necessitano di calcolatori e ci sono dei problemi di convergenza degli algoritmi. \newline 

\begin{tcolorbox}[colback=red!5!white,colframe=red!75!black] \label{q: 5}
"Necessitano di calcolatori", cioè ci saranno più conti da fare? 
\end{tcolorbox}

\newpage 

\subsection{Filtri analogici di Butterworth}

\footnote{ 
    Slide del prof | 9 - Progetto Filtri IIR 1 | pag 6 - 9\\
    Appunti | 2025-11-11 | pag 4 \\ 
    Appunti | 2025-11-12 | pag 1  
} 

\begin{tcolorbox}
    In commercio sono disponibili diversi filtri analogici. Nei seguenti capitoli ce ne saranno diversi a titolo di esempio, quelli che comunemente vengono utilizzati. \newline

    Inoltre, non dobbiamo ricordare tutte le formule dei filtri, ma il procedimento da specifiche agli argomenti da utilizzare, 
    perchè ormai le formule ed i software (nel nostro caso Matlab con le sue funzioni e filterDesigner) trovano i valori analitici di questi parametri. 
\end{tcolorbox}

Data la funzione del modulo del filtro analogico di Butterworth: 

{
    \Large 
    \begin{equation}
        \abs{H_c \left(\jmath \Omega\right)}^{2} 
        = 
        \dfrac{1}{1 + \left( \dfrac{\jmath \Omega}{\jmath \Omega_c}\right)^{2n}}
    \end{equation}
}

\begin{tcolorbox}[colback=red!5!white,colframe=red!75!black] \label{q: 6}
$\abs{H_c \left(\jmath \Omega\right)}$ al quadrato perchè si considera lo spettro di potenza per il teorema di Parseval? 
\end{tcolorbox}

In base al valore di n, avremo questi andamenti: 

\begin{figure}[h]
    \centering
    \includegraphics[scale = 1]{filtro di Butterworth diversi ordini di n.png}
\end{figure}  

Dalla seguente figura, si possono trarre le seguenti considerazioni: 

\begin{itemize}
    \item tutte le funzioni n-esime sono monotonicamente decrescenti 
    \item all'aumentare di n, il filtro si avvicina al filtro rettangolare ideale 
    \item per qualsiasi valore di n, alla frequenza di cutoff $\Omega_c$: 
    {
        \Large 
        \begin{equation}
            \abs{H_c \left(\Omega_c\right)}^{2} 
            = 
            \dfrac{1}{2} 
            \approx 
            -3 \text{ dB}
        \end{equation}
    } 
    \item è un filtro massimamente piatto 
\end{itemize}

\begin{tcolorbox}[colback=red!5!white,colframe=red!75!black] \label{q: 7}
Ha delle conseguenze e delle implicazioni sapere che è massimamente piatto?
\end{tcolorbox}

\newpage 

Ponendo: 

{
    \Large 
    \begin{equation}
        s = \jmath \Omega
    \end{equation}
}

avremo che la funzione $\abs{H \left(\jmath \Omega\right)}^{2}$ diventa: 

{
    \Large 
    \begin{equation}
         \abs{H_c \left(s\right)}^{2}_{s = \jmath \Omega}
            = 
            \left. 
            H_c \left(s\right)
            \cdot 
            H_c \left(- s\right)
            \right|_{s = \jmath \Omega}
            \\
            = 
            \left.
            \dfrac{1}{1 + \left( \dfrac{s}{\jmath \Omega_c}\right)^{2}}
            \right|_{s = \jmath \Omega}
    \end{equation}
}

Nel piano di Laplace, $H_c \left(s\right)$, in base al valore di n, avremo questi andamenti: 

\begin{figure}[h]
    \centering
    \includegraphics[scale = 0.5]{filtro di Butterworth diversi ordini di n piano Laplace.png}
\end{figure}  

dove: 

\begin{itemize}
    \item per n pari, i poli di $H_c \left(s\right)$ non si trovano sull'asse reale, a differenza dei filtri di Butterworth per n dispari 
    \item $\Omega_c$ è il raggio della circonferenza 
    \item per n dispari, la distanza di ogni polo è di $\dfrac{\pi}{n}$
\end{itemize}

\begin{tcolorbox}[colback=red!5!white,colframe=red!75!black] \label{q: 8}
Ha delle implicazioni avere dei poli nell'asse reale? La circonferenza $\Omega_c$ è di raggio unitario? 
\end{tcolorbox}

Come scritto precedentemente \ref{subsec: Specifiche "in frequenza"} \nameref{subsec: Specifiche "in frequenza"} a pag \pageref{subsec: Specifiche "in frequenza"}, 
passeremo dalle specifiche: 

\begin{figure}[h]
    \centering
    \includegraphics[scale = 1]{prototipo filtro passa basso.png}
\end{figure}  

alle seguente specifiche: 

\newpage 

\begin{figure}[h]
    \centering
    \includegraphics[scale = 0.5]{filtro di Butterworth specifiche}
\end{figure}  

dove, date le specifiche $\delta_1$, $\delta_2$, $\Omega_p$, $\Omega_s$, 
per la realizzazione di un filtro di Butterworth, sarà necessario trovare il valore di n che soddisfa le specifiche. \newline 

Di conseguenza, trovato il valore di n, possiamo trovare i coefficienti di $H_c \left(s\right)$. \newline 

Dall'espressione di $\abs{H_c \left(\jmath \Omega\right)}^{2}$:  

{
    \Large 
    \begin{equation}
        \abs{H_c \left(\jmath \Omega\right)}^{2} 
        = 
        \dfrac{1}{1 + \left( \dfrac{\jmath \Omega}{\jmath \Omega_c}\right)^{2n}}
    \end{equation}
}

si trova facilmente che le specifiche vengono soddisfatte se, ad un certo valore di n e $\Omega_c$ valgono contemporaneamente queste due relazioni: 

{
    \Large 
    \begin{equation}
        \dfrac{\Omega_p}{\Omega_c}
        = 
        \left(
            \dfrac{\delta_1}{1 - \delta_1}
        \right)^{\frac{1}{2n}}
    \end{equation}
}

{
    \Large 
    \begin{equation}
        \dfrac{\Omega_s}{\Omega_c}
        = 
        \left(
            \dfrac{1 - \delta_2}{\delta_2}
        \right)^{\frac{1}{2n}}
    \end{equation}
}

I valori di n e $\Omega_c$ vengono trovati per tentativi. \newline 


\begin{tcolorbox}
    Generalmente, in commercio, ci sono dei software che trovano i valori di n e $\Omega_c$ vengono trovati per tentativi. 
\end{tcolorbox}

I coefficienti di $H_c \left(s\right)$ si trovano analiticamente o da apposite tabelle. \newline 

\newpage 

\subsection{Filtri analogici di Chebyshev}

\footnote{ 
    Slide del prof | 9 - Progetto Filtri IIR 1 | pag 10\\ 
    Appunti | 2025-11-12 | pag 1  
} 

Basato sui polinomi di Chebyshev, al variare di n, possiamo costruire i filtri analogici di Chebyshev come si vede in figura: 

\begin{figure}[h]
    \centering
    \includegraphics[scale = 0.6]{polinomio di Chebyshev al variare di n.png}
\end{figure}  

dove: 

{
    \Large 
    \begin{equation}
        T_n \left(x\right) = \cos \left[ n \cos^{-1} \left(x\right)\right]
    \end{equation}
}

{
    \Large 
    \begin{equation}
        \begin{cases}
            T_0 \left(x\right) = 1 
            \\
            T_1 \left(x\right) = x 
            \\
            T_n \left(x\right) = 2x \cdot T_{n-1} \left(x\right) - T_{n-2} \left(x\right) \text{ per } n > 2 
        \end{cases}
    \end{equation}
}

\begin{tcolorbox}[colback=red!5!white,colframe=red!75!black] \label{q: 9}
L'uso dei polinomi di Chebyshev è quello di scrivere gli n-esimi coseni come una funzione polinomiale, giusto? Servono per l'interpolazione tra due punti? 
\end{tcolorbox}

\begin{tcolorbox}
    Se vuoi approfondire il tema sui polinomi di Chebyshev e perchè sono molto utili: 

    \begin{itemize}
        \item Chebyshev Polynomials by Physics and Math Lectures \\\url{https://www.youtube.com/watch?v=q8zGN2QwyU4&list=PL2uXHjNuf12aUiM6mgtng_mx07aB5o2zY}
        \item Quintic Equation From Chebyshev Polynomial by blackpenredpen\\\url{https://www.youtube.com/watch?v=Sgjy1uyYNbc} 
        \item Chebyshev methods: What are they and what can they do in practice? by QuantMinds TV\\\url{https://www.youtube.com/watch?v=FWp1X8m5XX4}
    \end{itemize}
\end{tcolorbox}

\newpage 

\subsubsection{Filtri analogici di Chebyshev di tipo I}

\footnote{ 
    Slide del prof | 9 - Progetto Filtri IIR 1 | pag 11 - 12\\ 
    Appunti | 2025-11-12 | pag 1  
} 

I filtri analogici di Chebyshev di tipo I hanno la seguente formula in modulo: 

{
    \Large 
    \begin{equation}
        \abs{H_c \left(\jmath \Omega\right)}^{2}
        = 
        \dfrac{1}{1 + \varepsilon^{2} T^{2}_{n} \left(\dfrac{\Omega}{\Omega_c}\right)}
    \end{equation}
}

dove: 

\begin{itemize}
    \item $T^{2}_{n} \left(\dfrac{\Omega}{\Omega_c}\right)$ è l'n-esimo polinomio di Chebyshev (visto nella pagina precedente) al quadrato 
    \item $\Omega_c$ è la frequenza di cut-off (che sarà una della specifica)
\end{itemize}


Due esempi di filtri di Chebyshev di tipo 1: 

\begin{figure}[h]
    \centering
    \includegraphics[scale = 0.3]{filtri di Chebyshev di tipo I esempi.png}
\end{figure}  

\begin{tcolorbox}[colback=red!5!white,colframe=red!75!black] \label{q: 10}
Perchè ci sono due figure di filtri di Chebyshev? Perchè uno inizia da 1 e uno da $1 - \delta_1$ ?
\end{tcolorbox}

Dalle seguenti figure possiamo trarre le seguenti considerazioni: 

\begin{itemize}
    \item nella banda passante, quindi in $\Omega < \Omega_c$, l'andamento del filtro non deve essere per forza essere sinusoidale come in figura; l'importante che il modulo del filtro si trovi tra 1 e $1 -\delta_1$
    \item ricordando che teoricamente, il filtro ideale dovrebbe essere rettangolare, nei filtri di Chebyshev di tipo I l'errore del modulo viene distribuito sulla banda passante (e questo giustifica perchè Chebyshev di tipo I è meglio di Butterworth) 
\end{itemize}

Date le specifiche $\delta_1$, $\delta_2$, $\Omega_c$ e $\Omega_s$ (dove generlmante $\delta_1$, $\delta_2$ sono dati in dB), 
bisogna calcolarsi gli oppurtuni valori di $\varepsilon$ e n di $\abs{H_c \left(\jmath \Omega\right)}^{2}$. \newline 

Dalla formula di $\abs{H_c \left(\jmath \Omega\right)}^{2}$ si deduce "facilmente" che, date le specifiche, avremo che i valori di $\varepsilon$ e n si trovino con queste formule: 

{
    \Large 
    \begin{equation}
        \varepsilon^{2}
        = 
        \dfrac{\delta_1}{1 - \delta_1}
    \end{equation}
}


{
    \Large 
    \begin{equation}
        n \cosh^{-1} \left(\dfrac{\Omega_s}{\Omega_c}\right)
        = 
        \cosh^{-1} \left[ 
            \sqrt{
                \dfrac{\left(1 - \delta_1\right) \left(1 - \delta_2\right) }{\delta_1 \delta_2}
            }
        \right]
    \end{equation}
}

da cui si possono ricavare i valori di $\varepsilon$ e n. \newline 

Proprio perchè il coseno iperbolico, cioè $\cosh$, p un calcolo complicato, generalmente i coefficienti della corrispondente funzione $H_c \left(s\right)$ si possono trovare 
tabulati per valori di $\delta_1 = 1$ e $\Omega_c = 1$. \newline 

\newpage 

\subsubsection{Filtri analogici di Chebyshev di tipo II}

\footnote{ 
    Slide del prof | 9 - Progetto Filtri IIR 1 | pag 13\\ 
    Appunti | 2025-11-12 | pag 1  
} 

I filtri analogici di Chebyshev di tipo II hanno la seguente formula in modulo: 

{
    \Large 
    \begin{equation}
        \abs{H_c \left(\jmath \Omega\right)}^{2}
        = 
        \dfrac{1}{1 + \dfrac{1}{\varepsilon^{2} T^{2}_{n} \left(\dfrac{\Omega_c}{\Omega}\right)}}
    \end{equation}
}

dove: 

\begin{itemize}
    \item $T^{2}_{n} \left(\dfrac{\Omega_c}{\Omega}\right)$ è l'n-esimo polinomio di Chebyshev (visto nella pagina precedente) al quadrato 
    \item $\Omega_c$ è la frequenza di cut-off (che sarà una della specifica)
\end{itemize} 

Un esempio di andamento del modulo di un filtro di Chebyshev di tipo II: 

\begin{figure}[h]
    \centering
    \includegraphics[scale = 0.5]{filtri di Chebyshev di tipo II esempio.png}
\end{figure}  

Essendo l'equiripple presente nella stop band, quindi per $\Omega > \Omega_s$, allora l'errore del filtro di Chebyshev di tipo II si trova in questa banda. \newline 

Per le specifiche, si adotta lo stesso approcio dei filtri di Tipo I di Chebyshev. \newline 

\newpage 

\subsection{Filtri analogici elittici}

\footnote{ 
    Slide del prof | 9 - Progetto Filtri IIR 1 | pag 13 - \\ 
    Appunti | 2025-11-12 | pag 1  
} 